% !TeX root = ../教材题库_全部条目_转写工作稿.tex
% Source: XJ-B1-C03-SU-P089-Y1112
% 第3章 / 小结与复习 小结与复习 / 教材 p.89 / PDF 第96页
% OCR未校对：公式与文字均可能有误；未完成数学TeX转写；禁止直接发给学生
% 原题及相邻原解截图（核对用）：../题目原图/XJ-B1-C03-SU-P089-Y1112-01.webp
% 本文件只在转写完成后作为题干使用；OCR 条目尚不是可靠题干。
\begin{bankitem}{第3章 小结与复习 第3题}{第3章 / 小结与复习 小结与复习 / 教材 p.89 / PDF 第96页}{XJ-B1-C03-SU-P089-Y1112}{OCR未校对：公式与文字均可能有误；未完成数学TeX转写；禁止直接发给学生}
\begin{ocrdraft}
[教材p.89 / PDF第96页]
\par
3．根据函数的解析式或图象来探索它的性质，是数学的重要课题．函数的变
\par
化无非是变大变小，函数的单调性描述了函数的变化过程和趋势，是函数最重要
\par
的特征．如何判断一个函数的单调性？更进一步，你能用数学语言来表述吗？函
\par
数最值的几何意义如何理解？
\end{ocrdraft}
\end{bankitem}
